\documentclass[12pt,a4paper]{article}

\usepackage[margin=1in]{geometry}
\usepackage[hidelinks]{hyperref}
\usepackage{setspace}
\usepackage[ngerman]{babel}
\usepackage{changepage}
\usepackage{newtxtext}

\setlength{\parskip}{0.5\baselineskip}
\setlength{\parindent}{1.5em}

\onehalfspacing
\pagestyle{empty}

\begin{document}

Guten Tag allerseits,\\%
Ich hoffe es geht Ihnen allen gut. Wir stehen schon in der Mitte des Quartals und ebenso befinden wir uns auch nur einige Schritten an der Mitte des Romans vorbei. Heute werden Lila und ich über Kapitel fünf, das in der Stadt Heidelberg stattfindet, sprechen.\par
%
Einiges zu Heidelberg:\\%
Die Stadt liegt in Baden-Württemberg an einem Fluss namens der Neckar. Die Neckarauen die oft in diesem Kapitel erwähnt werden, sind die Gebiete direkt nebenan dem Fluss. Im Englischen heißt „die Aue“ “the Floodplain”. Der Protagonist redet auch von den alten Gebäuden. In Heidelberg sind diese richtig alt: Heidelberg so wie wir die Stadt heute kennen gibt es seit dem zwölften Jahrhundert, aber Menschen hatten bereits in den Fünften dort gewohnt. Die Universität in Heidelberg ist am Ende des vierzehnten Jahrhunderts gegründet worden, ist die älteste Universität Deutschlands und gilt immernoch als einen wichtigen akademischen und kulturellen Zentrum.\par
%
\begin{center}
    \textit{Lila spricht über Heidelberg\\ WICHTIG–SPRICHT ÜBER HOTEL ALT HEIDELBERG UND BOMBEN IN HEIDELBERG}%
\end{center}
\par
Nun, was passiert eigentlich in dem Heidelberg-Kapitel mit dem namenlosen Protagonist? Ich möchte gerne auf einer Stelle im Text hinweisen. Folgendes steht bei mir auf Textseite 82:\\%
\vspace{-2.5em}
\begin{center}
\textit{Na, jedenfalls mache ich mit dem Schlüssel die Hoteltür auf, ich bin ja immer noch im Hotel Alt Heidelberg, während ich mich an das blöde Spiel auf Madeira erinnere, und dann gehe ich ins Zimmer und mache erst mal das Licht an, weil es ziemlich düster ist da drinnen.“}\\%
\end{center}%
\vspace{-.5em}
Der Erzähler sagt das gleich nach Einem seiner üblichen Flashbacks. Der Schlüssel düftet „sauber, aber irgendwie auch ungewaschen“, ähnlich wie einer Seife die er als Kind in der Hosentasche in Madeira dabei gehabt hatte. Er spricht von der Zeitschrift „Bravo“ und wie diese ihm weggenommen wurde. Interessant an dieser Stelle ist, als er sagt er sei ja noch immer im Hotel Alt Heidelberg. Es ist fast als ob er sich aktiv daran erinnern muss, dass er kein Kind im Urlaub mehr ist. Dabei erinnert er sich so genau an dieser Zeitschrift, dass er weiß welche Serie die erste Person mit weißen Bermudashorts und Collegeschuhe beinhaltet. Wir haben hier den sogenannten „Markenfetischismus“ der so oft im Roman auftaucht. Hier aber in wahnsinnigen Grad. Er schimpft die Bravo dadurch, dass er die enthaltene Artikeln als schamhaft bezeichnet und die Zeitschrift im Ganzen als gottverdammt, jedoch ist er scheinbar unvergnügt, als er sie am Ende des Flashbacks nicht wiederbekommt. Wenn er sie wirklich so geringgeschätzt hätte, da würde er sich freuen sie nicht wieder zu bekommen. Da steckt aber eine gewisse Ambivalenz drin, dass er einfach ergänzt er hätte sie nicht wiedererhalten. Man könnte den Markenfetischismus als eine „Pseudogemeinschaft“ bezeichnen: eine Gemeinschaft, weil er eine Weise, Sozialverhältnisse zu deuten bietet, aber auch pseudo, also falsch, weil diese Weise nicht immer stimmt und komplett küntslich geschaffen ist. Vielleicht ist die Ambivalenz an dieser Stelle in gewissermaßen eine Anerkennung dessen? Doch das ist eigentlich sehr spekulativ. Hier, wie fast ständig im Buch, weiß der Leser nicht genau was mit irgendeiner Aussage gemeint ist.\par
%
Ich möchte aber zu meinem originallen Zitat zurück. Der Erzähler nennt das Spiel blöd, aber wie wir schon besprochen haben, unterdrückt er oft seine Gefühle. Die Sekundärliteratur spricht oft davon, dass Ich im Buch keinen Zweck hat. Er reist durch Deutschland, will zwar nach Karlsruhe weil er noch nie da war, trotzdem ändert er sich ohne es tiefer zu überlegen den Zielort, nur weil er nicht mit Horx sprechen will. Natürlich veranlasst das Ich das Gespräch mit Horx, daher ist das Ichs Benehmen widersprüchlich, aber das will ich jetzt nicht besprechen. Die Hauptsache an dieser Erinnerung ist, meines Erachtens, der Satz:\\%
\vspace{-2.5em}
\begin{center}
    \textit{Die anderen waren alle erwachsen und haben gelacht, weil es ihnen auf eine andere Art Spaß gemacht hat als mir, aber ich, ich bin vollkommen zielstrebig auf den Pool zugelaufen.}\\%
\end{center}
\vspace{-.5em}
Der Preis bedeutet für ihn nichts, und er interessiert sich überhaupt nicht für den Ort an dem er ist, aber bloß das Ziel an sich hat für ihn Wert getragen. Hier ist es interessant die Parallelen zu merken, er ist alleine in einer Art Urlaub bei einem Hotel und er weiß anfangs nicht genau, was er tun will. Natürlich wird die Erinnerung durch den Duft des Schlüssels ausgelöst; es wird ja oft gesagt der Geruchssinn sei von der fünf Sinne am stärksten mit der Erinnerung verbunden. Wichtiger ist es aber zu überlegen, was genau diese Situationen unterschieden. Ich habe es bereits gesagt: die Anwesenheit eines Ziels. Er muss sich so zu sagen zusammenreißen in der Gegenwart zu sein, wie oft der Fall ist, da er eigentlich danach sehnt, von jemandem anderen ein Ziel gegeben zu werden, so wie damals der Hotelangestellter, trotz aller Leiden, wie Blut und Wunde am Bein. Das ist auch warum er gegenwärtig den Hotelmann fragt, wohin die junge Leute nachts gehen, und direkt zu Maxs Bar hinfährt und dann später auch mit Eugen auf die Party geht.\par
%
Meine These ist, dass er das Spiel nicht wirklich Blöd fand, sondern danach sehnt.
\end{document}
